\documentclass[10pt]{amsart}
\usepackage[margin=1in]{geometry}
\usepackage[T1]{fontenc}
\usepackage{lmodern}
\usepackage{amsmath,amssymb}
\usepackage[hidelinks]{hyperref}
\usepackage{enumitem}
\setlist[itemize]{leftmargin=1.4em,itemsep=3pt}
\title{Verification guide for stabilizer-state identification}
\date{15 October 2026}
\begin{document}
\maketitle
\noindent
This repository contains two independent Lean 4 projects accompanying papers
by Masahito Hayashi and Yimin Lu. This guide describes the accepted formal
scope and how to verify an exact release. It does not certify a future commit.

\section{Paper correspondence}
\subsection*{arXiv:2609.13923}
\emph{Partial Stabilizer Learning under Fixed Commuting Constraints and the
Absence of a Copy-Rate Discount}~\cite{partial}.
The root project uses the namespace \texttt{StabilizerLearning}.
Its main endpoint is
\begin{center}\small
\texttt{StabilizerLearning.fixedW\_noCopyRateDiscount}.
\end{center}
Its exact type is \texttt{StabilizerLearning.FixedWNoCopyRateDiscount}.
The complete original fixed-$W$ target includes the common-$J$ converse,
four uniform limits to zero, four uniform limits to one under the original
condition $k_n-n\to+\infty$, and the exact copy-rate identities.
The formalization retains every prime local dimension, the actual states
and POVMs, and the original order of worst-case optimization.
This is closure of the original main theorem, not a claim that every
auxiliary statement in the manuscript is formalized.
The itemized coverage and remaining interfaces are in
\texttt{docs/theorem-map.md}.

\subsection*{arXiv:2610.00266}
\emph{Exact Critical Curve for Uniform Stabilizer-State
Identification}~\cite{critical}.
The independent project under \texttt{proofs/arxiv-2610.00266/}
uses the namespace \texttt{ExactCriticalCurve}.
Its three main endpoints are
\texttt{finiteSpectrum}, \texttt{criticalCurve} and
\texttt{criticalCurveTails}, all in that namespace.
They cover the actual finite Gram spectrum and optimal PGM for $n,k\geq1$,
the literal limiting critical curve for every fixed integer offset $c$,
the exact first-moment identity, and both coefficient-one tails.
Additional results cover actual stochastic-group orders, ordinary/nonzero/zero
spectral-mass limits, and the exceptional-amplitude inequality.
The nested \texttt{docs/theorem-map.md} gives the precise quantifiers and
endpoint conventions.

The package does not certify the displayed numerical evaluation at $c=0$,
a joint limit with a growing offset, full spectral-measure weak
convergence, or a complete exceptional shortening equality merely from
the amplitude bound. The formal convergence proof uses sufficient summable
bounds and does not certify the paper's exact displayed auxiliary majorant.

\section{Meaning of verification}
The strict gates check compilation, recursive axiom dependencies,
independently expanded theorem types, and required original proof bodies.
Only \texttt{propext}, \texttt{Classical.choice} and \texttt{Quot.sound}
are permitted by the recursive audits; placeholders or new mathematical
axioms do not pass.

These checks certify formal declarations relative to the pinned Lean kernel
and libraries. They do not independently decide whether every formal
definition faithfully expresses the intended paper semantics.
Recorded semantic reviews are AI reviews.
Independent human mathematical review remains outstanding.
\newpage

\section{Reproducing an immutable release}
Supported pinned compiler platforms are Linux x86\_64 and macOS ARM64.
Install Git, Python 3.10 or newer, curl, tar and zstd.
Keep the two projects' dependencies, compiler directories and build outputs
separate. The second complete audit may take several hours.

Obtain the repository from
\href{https://github.com/Imine20/Lean-formalization-for-stabilizer-state-identification}
{the release repository}, then check out the exact commit to be tested.
Run all commands below from its root.
Record \texttt{git rev-parse HEAD} and \texttt{git rev-parse HEAD\^{}\{tree\}}
and confirm that \texttt{git status --porcelain} is empty.

\subsection*{First paper}
\begin{verbatim}
python3 scripts/bootstrap.py
./scripts/lake.sh exe cache get $(cat scripts/cache-imports.txt)
python3 scripts/check.py --require-main-proof
\end{verbatim}
Also kernel-check the independently expanded original Main statement:
\begin{verbatim}
expected=docs/validation/2026-10-11-direct-integration
./scripts/lake.sh env lean \
  "$expected/independent-expected/IndependentAcceptance.lean"
\end{verbatim}
Require successful command exits and the independent acceptance marker.
The strict gate writes its source hashes and logs under
\texttt{artifacts/check/}.

\subsection*{Second paper}
\begin{verbatim}
python3 scripts/verify_arxiv_2610.py \
  --dest /absolute/new/reproduction/path
\end{verbatim}
The destination must not already exist.
The wrapper reconstructs the immutable checked-out commit, obtains its pinned
public dependencies, begins with no project proof artifacts, and runs the
complete strict release gate.
Official Mathlib dependency caches are used by default.
The option \texttt{--dependency-mode source} instead compiles the required
dependency modules from their pinned source.
Neither mode rebuilds the Lean compiler from C++ sources.

\section{Pins, provenance and release acceptance}
Both projects pin Lean 4.24.0 and Mathlib at the same commits:
\begin{center}\footnotesize\ttfamily
Lean: 797c613eb9b6d4ec95db23e3e00af9ac6657f24b\\
Mathlib: f897ebcf72cd16f89ab4577d0c826cd14afaafc7
\end{center}
Each project retains its own complete transitive lock.
Historical evidence retains its original frozen inputs and execution paths.
Those paths are records of earlier checks, not public reproduction dependencies.

GitHub Actions runs both strict gates on pushes and pull requests.
For the final \texttt{main} commit, require both jobs to succeed and inspect
the artifacts for the exact commit/tree, unchanged input hashes,
all command exits and complete logs.
A prior successful run is historical evidence, not acceptance of later edits.
The release preparation preserves the mathematical source, paper snapshots,
version pins and verification tooling.

\section{Citation for manuscript revisions}
Cite the versioned
\href{https://github.com/Imine20/Lean-formalization-for-stabilizer-state-identification/releases/tag/v1.0.0}
{v1.0.0 release}, using \texttt{CITATION.bib}, together with the relevant
paper or papers. The release record identifies the tested commit and tree;
the evidence bundle retains the new command logs, axiom audits and type/body
checks for both projects. A source-file checksum manifest binds the cited
snapshot. The moving \texttt{main} branch is not a substitute for the
version tag. The documented scope boundaries continue to apply.

\begin{thebibliography}{9}
\bibitem{partial}
M.~Hayashi and Y.~Lu,
\emph{Partial Stabilizer Learning under Fixed Commuting Constraints and the
Absence of a Copy-Rate Discount},
\href{https://arxiv.org/abs/2609.13923}{arXiv:2609.13923v1} (2026).
\bibitem{critical}
M.~Hayashi and Y.~Lu,
\emph{Exact Critical Curve for Uniform Stabilizer-State Identification},
\href{https://arxiv.org/abs/2610.00266}{arXiv:2610.00266v1} (2026).
\end{thebibliography}
\end{document}
